% Generated by scripts/convert_report.py from source/report.md.


\FloatBarrier

\Needspace{6\baselineskip}

\section{GIỚI THIỆU ĐỀ TÀI VÀ CÔNG NGHỆ SỬ DỤNG}


\subsection{Giới thiệu đề tài}


\subsubsection{Bối cảnh và động lực thực hiện}


Trong các ứng dụng học máy, việc xây dựng một mô hình có độ chính xác tốt mới chỉ là một phần của quá trình phát triển. Khi đưa mô hình vào sử dụng thực tế, hệ thống còn phải đáp ứng những yêu cầu về quản lý phiên bản, tự động hóa triển khai, giám sát hiệu năng và xử lý sự cố. Nếu các công đoạn này thiếu sự liên kết, việc cập nhật mô hình có thể gây ra sai lệch phiên bản, suy giảm chất lượng dịch vụ hoặc khó khăn khi cần khôi phục hệ thống.


MLOps (Machine Learning Operations) giải quyết những vấn đề trên bằng cách kết hợp Machine Learning với các phương pháp DevOps, hướng đến việc quản lý và tự động hóa toàn bộ vòng đời mô hình, từ huấn luyện, đánh giá đến triển khai và vận hành.


Xuất phát từ nhu cầu đó, đề tài \textbf{xây dựng nền tảng MLOps trên Amazon Web Services (AWS)} được thực hiện nhằm nghiên cứu và triển khai một quy trình học máy hoàn chỉnh trên môi trường cloud, đồng thời kiểm chứng khả năng quản lý phiên bản và phát hành mô hình an toàn.


\subsubsection{Mục tiêu của đề tài}


Mục tiêu chính là xây dựng nền tảng MLOps trên AWS, tích hợp các thành phần huấn luyện, quản lý mô hình, triển khai dịch vụ và giám sát thành một quy trình thống nhất.


Các mục tiêu cụ thể bao gồm:


\begin{itemize}

\item 

\textbf{Quản lý hạ tầng:} Sử dụng Infrastructure as Code để triển khai môi trường Kubernetes và phân bổ tài nguyên tính toán phù hợp cho từng workload.


\item 

\textbf{Tự động hóa huấn luyện:} Thực hiện training trên cloud, theo dõi thí nghiệm, đánh giá chất lượng và quản lý phiên bản mô hình.


\item 

\textbf{Triển khai dịch vụ dự đoán:} Đóng gói mô hình thành API và xây dựng quy trình CI/CD kết hợp GitOps.


\item 

\textbf{Giám sát và kiểm soát phát hành:} Theo dõi hiệu năng, triển khai canary và hỗ trợ khôi phục phiên bản ổn định khi bản phát hành không đạt yêu cầu.


\end{itemize}

\subsubsection{Phạm vi thực hiện}


Đề tài sử dụng bài toán dự đoán thời lượng chuyến đi taxi tại New York làm ứng dụng minh họa. Phạm vi triển khai được chia thành bốn giai đoạn:


\begin{enumerate}

\item 

\textbf{AWS Infrastructure:} Xây dựng hạ tầng bằng Terraform, triển khai Amazon EKS và cơ chế quản lý tài nguyên Kubernetes.


\item 

\textbf{Training Pipeline \& Model Registry:} Huấn luyện mô hình trên cloud, đánh giá chất lượng và quản lý phiên bản bằng MLflow.


\item 

\textbf{Model Serving \& CI/CD:} Xây dựng FastAPI, đóng gói container và tổ chức các workflow kiểm thử, huấn luyện, phát hành.


\item 

\textbf{Progressive Delivery \& Observability:} Triển khai canary, giám sát hiệu năng và kiểm chứng cơ chế rollback.


\end{enumerate}

Hệ thống được triển khai trong môi trường phát triển và thực nghiệm. Các chức năng như tự động phát hiện data drift, tái huấn luyện liên tục, High Availability và Disaster Recovery toàn diện chưa thuộc phạm vi hiện thực.


\subsection{Bài toán ứng dụng: Dự đoán thời lượng chuyến đi taxi}


\subsubsection{Mô tả bài toán}


Bài toán \textbf{NYC Taxi Trip Duration Prediction} nhằm dự đoán thời lượng của một chuyến taxi dựa trên khu vực đón khách, khu vực trả khách và thời điểm bắt đầu chuyến đi. Đây là bài toán hồi quy (Regression), với đầu ra là thời gian di chuyển dự kiến tính bằng phút.


Mô hình chỉ sử dụng những thông tin có thể biết trước khi chuyến đi diễn ra. Các thuộc tính phát sinh sau khi kết thúc chuyến đi, chẳng hạn thời điểm trả khách hoặc giá cước, không được sử dụng làm đặc trưng đầu vào nhằm tránh hiện tượng rò rỉ dữ liệu (Data Leakage).


Trong đề tài, bài toán đóng vai trò ứng dụng thực tế để kiểm chứng quy trình MLOps, thay vì tập trung phát triển một hệ thống dự đoán giao thông chuyên biệt.


\subsubsection{Dữ liệu và quá trình chuẩn bị}


Đề tài sử dụng bộ dữ liệu \textbf{Yellow Taxi Trip Records} do New York City Taxi and Limousine Commission (NYC TLC) công bố. Dữ liệu gồm các bản ghi chuyến đi với thông tin thời gian, khu vực đón/trả khách và những thuộc tính liên quan.


Sau khi loại bỏ các bản ghi không hợp lệ, hệ thống lấy mẫu xác định với seed bằng 42, tối đa 100.000 bản ghi mỗi tháng. Dữ liệu được chia theo thứ tự thời gian:


\begin{table}[!htbp]\centering
\fontsize{12}{15}\selectfont\setstretch{1.05}
\setlength{\tabcolsep}{4pt}\setlength{\arrayrulewidth}{0.4pt}\renewcommand{\arraystretch}{1.2}
\setlength{\ReportTableWidth}{\dimexpr\linewidth-4\arrayrulewidth\relax}
\begin{tabular}{|L{\dimexpr 0.34\ReportTableWidth-2\tabcolsep\relax}|L{\dimexpr 0.34\ReportTableWidth-2\tabcolsep\relax}|L{\dimexpr 0.32\ReportTableWidth-2\tabcolsep\relax}|}
\hline
\textbf{Tập dữ liệu} & \textbf{Khoảng thời gian} & \textbf{Số bản ghi} \\
\hline
Training & 01–02/2025 & 200.000 \\\hline
Validation & 03/2025 & 100.000 \\\hline
Testing & 04/2025 & 100.000 \\\hline
\textbf{Tổng cộng} & \textbf{01–04/2025} & \textbf{400.000} \\\hline
\end{tabular}
\caption{Phân chia dữ liệu theo thời gian}\label{tab:01}
\end{table}

Cách chia này sử dụng dữ liệu quá khứ để huấn luyện và các giai đoạn tiếp theo để đánh giá, hạn chế việc trộn dữ liệu tương lai vào tập huấn luyện.


Thời lượng chuyến đi được tính từ chênh lệch giữa thời điểm trả khách và đón khách, giới hạn trong khoảng 1–180 phút. Mô hình sử dụng năm đặc trưng:


\begin{table}[!htbp]\centering
\fontsize{12}{15}\selectfont\setstretch{1.05}
\setlength{\tabcolsep}{4pt}\setlength{\arrayrulewidth}{0.4pt}\renewcommand{\arraystretch}{1.2}
\setlength{\ReportTableWidth}{\dimexpr\linewidth-3\arrayrulewidth\relax}
\begin{tabular}{|L{\dimexpr 0.38\ReportTableWidth-2\tabcolsep\relax}|L{\dimexpr 0.62\ReportTableWidth-2\tabcolsep\relax}|}
\hline
\textbf{Đặc trưng} & \textbf{Ý nghĩa} \\
\hline
\code{PULocationID} & Mã khu vực đón khách \\\hline
\code{DOLocationID} & Mã khu vực trả khách \\\hline
\code{pickup_hour} & Giờ bắt đầu chuyến đi \\\hline
\code{pickup_day_of_week} & Thứ trong tuần \\\hline
\code{pickup_month} & Tháng bắt đầu chuyến đi \\\hline
\end{tabular}
\caption{Các đặc trưng đầu vào của mô hình}\label{tab:02}
\end{table}

Hai mã khu vực được xử lý như đặc trưng phân loại, trong khi các đặc trưng thời gian được trích xuất từ thời điểm đón khách. Dữ liệu được lưu trữ trên Amazon S3 để phục vụ quá trình huấn luyện trên cloud.


\subsubsection{Mô hình và phương pháp đánh giá}


Đề tài sử dụng \textbf{LightGBM (Light Gradient Boosting Machine)}, một thuật toán Gradient Boosted Decision Trees phù hợp với dữ liệu dạng bảng và hỗ trợ đặc trưng phân loại.


Để đánh giá chất lượng, hệ thống sử dụng mô hình baseline dự đoán theo trung vị thời lượng chuyến đi của tập huấn luyện. Kết quả LightGBM được đánh giá bằng hai chỉ số:


\begin{itemize}

\item 

\textbf{MAE (Mean Absolute Error):} Sai số tuyệt đối trung bình giữa giá trị dự đoán và thực tế, được sử dụng làm chỉ số đánh giá chính.


\item 

\textbf{RMSE (Root Mean Squared Error):} Căn bậc hai của sai số bình phương trung bình, giúp phản ánh rõ hơn các trường hợp có sai số lớn.


\end{itemize}

Quá trình lựa chọn mô hình dựa trên tập validation, trong khi tập test được sử dụng để đánh giá sau kiểm duyệt. Kết quả validation cũng là cơ sở cho \textbf{Validation Gate} quyết định một mô hình ứng viên có đủ điều kiện đăng ký phiên bản mới trong Model Registry hay không.


Quy trình huấn luyện được trình bày tại Chương 3 và kết quả đánh giá được phân tích tại Chương 4.


\subsection{Công nghệ sử dụng và quyết định thiết kế}


Các công nghệ được lựa chọn theo vai trò trong vòng đời MLOps, kết hợp dịch vụ AWS với các công cụ mã nguồn mở nhằm đáp ứng yêu cầu triển khai, tự động hóa và giám sát hệ thống.


\subsubsection{Nhóm hạ tầng và điều phối tài nguyên}


\begin{table}[!htbp]\centering
\fontsize{12}{15}\selectfont\setstretch{1.05}
\setlength{\tabcolsep}{4pt}\setlength{\arrayrulewidth}{0.4pt}\renewcommand{\arraystretch}{1.2}
\setlength{\ReportTableWidth}{\dimexpr\linewidth-3\arrayrulewidth\relax}
\begin{tabular}{|L{\dimexpr 0.35\ReportTableWidth-2\tabcolsep\relax}|L{\dimexpr 0.65\ReportTableWidth-2\tabcolsep\relax}|}
\hline
\textbf{Công nghệ} & \textbf{Vai trò} \\
\hline
\textbf{AWS} & Cung cấp môi trường cloud và các dịch vụ tính toán, mạng, lưu trữ \\\hline
\textbf{Terraform} & Quản lý hạ tầng theo mô hình Infrastructure as Code \\\hline
\textbf{Amazon EKS / Kubernetes} & Điều phối container và các workload của hệ thống \\\hline
\textbf{Karpenter} & Tự động cấp phát và thu hồi EC2 node theo nhu cầu \\\hline
\textbf{Amazon S3 / ECR / EBS} & Lưu trữ datasets, model artifacts, container images và dữ liệu bền vững \\\hline
\end{tabular}
\caption{Công nghệ hạ tầng và điều phối tài nguyên}\label{tab:03}
\end{table}

Terraform đảm nhiệm việc quản lý tài nguyên AWS, trong khi Kubernetes tổ chức và vận hành các thành phần ứng dụng.


Hệ thống sử dụng \textbf{EC2 Spot cho training workloads} nhằm tối ưu chi phí đối với các tác vụ có thời gian thực thi hữu hạn, đồng thời sử dụng \textbf{On-Demand cho model serving} để hạn chế gián đoạn do thu hồi Spot.


\subsubsection{Nhóm huấn luyện, quản lý mô hình và serving}


\begin{table}[!htbp]\centering
\fontsize{12}{15}\selectfont\setstretch{1.05}
\setlength{\tabcolsep}{4pt}\setlength{\arrayrulewidth}{0.4pt}\renewcommand{\arraystretch}{1.2}
\setlength{\ReportTableWidth}{\dimexpr\linewidth-3\arrayrulewidth\relax}
\begin{tabular}{|L{\dimexpr 0.35\ReportTableWidth-2\tabcolsep\relax}|L{\dimexpr 0.65\ReportTableWidth-2\tabcolsep\relax}|}
\hline
\textbf{Công nghệ} & \textbf{Vai trò} \\
\hline
\textbf{Python / LightGBM} & Tiền xử lý dữ liệu và huấn luyện mô hình \\\hline
\textbf{MLflow Tracking} & Theo dõi tham số, metrics và artifacts của các thí nghiệm \\\hline
\textbf{MLflow Model Registry} & Quản lý phiên bản, trạng thái phê duyệt và alias \code{champion} \\\hline
\textbf{PostgreSQL} & Lưu trữ metadata của MLflow \\\hline
\textbf{FastAPI} & Cung cấp HTTP API cho dịch vụ dự đoán \\\hline
\end{tabular}
\caption{Công nghệ huấn luyện, quản lý mô hình và serving}\label{tab:04}
\end{table}

MLflow giúp liên kết mỗi phiên bản mô hình với kết quả huấn luyện và thông tin nguồn, hỗ trợ việc truy vết và so sánh các thí nghiệm.


Mô hình được phê duyệt sẽ được đóng gói cùng ứng dụng FastAPI thành container image, bảo đảm phiên bản phục vụ gắn với một bản phát hành xác định.


\subsubsection{Nhóm CI/CD, triển khai và giám sát}


\begin{table}[!htbp]\centering
\fontsize{12}{15}\selectfont\setstretch{1.05}
\setlength{\tabcolsep}{4pt}\setlength{\arrayrulewidth}{0.4pt}\renewcommand{\arraystretch}{1.2}
\setlength{\ReportTableWidth}{\dimexpr\linewidth-3\arrayrulewidth\relax}
\begin{tabular}{|L{\dimexpr 0.35\ReportTableWidth-2\tabcolsep\relax}|L{\dimexpr 0.65\ReportTableWidth-2\tabcolsep\relax}|}
\hline
\textbf{Công nghệ} & \textbf{Vai trò} \\
\hline
\textbf{GitHub Actions} & Tự động kiểm thử, huấn luyện và phát hành \\\hline
\textbf{AWS IAM / OIDC} & Xác thực và phân quyền cho các workflow truy cập AWS \\\hline
\textbf{Argo CD} & Đồng bộ ứng dụng trên Kubernetes theo cấu hình GitOps \\\hline
\textbf{Argo Rollouts} & Triển khai canary và kiểm soát quá trình phát hành \\\hline
\textbf{AWS Application Load Balancer} & Phân phối lưu lượng giữa stable và canary \\\hline
\textbf{Prometheus / Grafana} & Thu thập, truy vấn và trực quan hóa metrics \\\hline
\textbf{k6} & Tạo tải HTTP và đánh giá hiệu năng dịch vụ \\\hline
\end{tabular}
\caption{Công nghệ CI/CD, triển khai và giám sát}\label{tab:05}
\end{table}

GitHub Actions đảm nhiệm các quy trình CI/CD, còn Argo CD đồng bộ trạng thái hệ thống theo cấu hình được quản lý trong Git. Argo Rollouts kết hợp ALB và Prometheus để triển khai canary và đánh giá phiên bản mới trước khi chuyển toàn bộ lưu lượng.


Một quyết định thiết kế quan trọng là \textbf{tách biệt kiểm duyệt chất lượng mô hình và kiểm duyệt bản phát hành}. Mô hình được Validation Gate phê duyệt chưa đồng nghĩa có thể đưa ngay vào phục vụ; phiên bản mới còn phải đáp ứng các điều kiện vận hành trong quá trình triển khai.


Cách tổ chức này giúp nền tảng kiểm soát cả chất lượng dự đoán và rủi ro phát hành mô hình.

